\section{Conclusion}\label{sec:conclusion}

This project asked whether an attention-based multivariate Bi-LSTM combined with volatility-aware risk management could make a leveraged Bitcoin CFD trader more stable and better at preserving capital. The answer is yes for capital preservation and no for prediction. The Alpha Agent could not forecast the next day's direction better than a coin flip, and neither could any simple rule I tested. More features (H1) and attention (H2) gave at best inconsistent gains in forecast error. The volatility-aware Risk Engine (H3) cut maximum drawdown by 12.6 percentage points in 2025 and 14.1 points on unseen 2026 data compared with fixed rules, significantly both times and at every leverage cap, without lowering the Sharpe ratio. It did not reduce liquidations, and at 20x leverage its volatility stops caused two.

For a retail trader the practical message is that, on this evidence, deep learning was useful for judging how risky tomorrow is and not for judging which way the price will go. Because the framework is decoupled, the two findings do not depend on each other. A better direction model could be swapped in without touching the risk side.

Several pieces of future work follow directly from the results. The stop distance should be capped inside the liquidation distance. The volatility forecast should be recalibrated on validation data to remove its upward bias, in the same way that Platt scaling adjusts the jump probabilities. The backtest should include slippage and overnight financing costs. For direction, training a classifier instead of an MSE regressor, forecasting a weekly horizon where momentum effects have been found \cite{liu2021}, or adding derivatives and order-book data are more promising than a larger network. Any of these would need new test data, since 2025 and 2026 have now both been looked at, and the same tests should be repeated on other assets such as Ethereum.
